% ============================================================
%  sections/01_introduction.tex
%  1. Giới thiệu  (~0.5 trang)
% ============================================================

\section{Giới Thiệu}
\label{sec:intro}

% --- Bối cảnh ---
Sự bùng nổ của các mô hình ngôn ngữ lớn (LLM) như GPT-4, LLaMA và
Qwen trong những năm gần đây đã tạo ra khả năng tạo văn bản tự động
với chất lượng gần tương đương con người~\cite{brown2020gpt3,
touvron2023llama}. Trên các nền tảng mạng xã hội (MXH) đa ngôn ngữ,
điều này tạo ra rủi ro nghiêm trọng: nội dung do máy tạo ra (Machine-Generated
Text --- MGT) có thể được phát tán quy mô lớn nhằm mục đích gây nhiễu
thông tin, thao túng dư luận, hoặc tự động hóa spam~\cite{zellers2019grover}.

% --- Khoảng trống nghiên cứu ---
Hầu hết các bộ phát hiện MGT hiện nay được xây dựng và đánh giá chủ
yếu trên tiếng Anh~\cite{solaiman2019gpt2detector,mitchell2023detectgpt}.
Khi áp dụng sang ngôn ngữ khác, hiệu năng sụt giảm nghiêm trọng do
thiên lệch ngôn ngữ trong quá trình huấn luyện. Ngoài ra, nhiều phương
pháp tiên tiến đòi hỏi GPU A100 (40--80 GB VRAM), xa tầm tay của hầu
hết nhà nghiên cứu và tổ chức tại các nước đang phát triển.

% --- Đóng góp ---
Bài báo này đề xuất và đánh giá ba phương pháp phát hiện MGT trên
dữ liệu MXH đa ngôn ngữ, với các đóng góp chính sau:

\begin{enumerate}
    \item \textbf{Bộ dữ liệu đa ngôn ngữ:} Xây dựng MultiSocial-4L, bộ
          dữ liệu gồm 15.985 bài đăng MXH cân bằng (human/machine) trên bốn
          ngôn ngữ EN, VI, ZH, AR, được tạo từ các corpus thực tế và các LLM
          hiện đại.

    \item \textbf{Đánh giá hệ thống:} So sánh toàn diện bốn phương pháp ---
          zero-shot với bộ phát hiện tiền huấn luyện (RoBERTa), tinh chỉnh
          mô hình đa ngôn ngữ (XLM-RoBERTa), phát hiện thống kê dựa trên
          Perplexity (Qwen2.5-1.5B), và học chuyển giao đa ngôn ngữ
          (leave-one-language-out) --- trên cùng một bộ dữ liệu.

    \item \textbf{Khả thi trên phần cứng phổ thông:} Toàn bộ pipeline
          thực nghiệm được thiết kế và chạy thành công trên Google Colab T4
          (16 GB VRAM), với tổng chi phí API ước tính \$1,50.
\end{enumerate}

% --- Cấu trúc bài báo ---
Phần còn lại của bài báo được tổ chức như sau: Phần~\ref{sec:related}
điểm lại các công trình liên quan; Phần~\ref{sec:dataset} mô tả quá
trình xây dựng bộ dữ liệu; Phần~\ref{sec:method} trình bày bốn phương
pháp thực nghiệm; Phần~\ref{sec:results} báo cáo và phân tích kết quả;
Phần~\ref{sec:conclusion} tổng kết và đề xuất hướng nghiên cứu tiếp theo.
